\subsection{设计目标与软件边界}

RM26赛季空中机器人的软件由机载计算平台、飞控和电控板三部分协同完成。机载计算平台负责视觉识别、定位导航与飞行指令生成，飞控负责姿态稳定和旋翼执行，电控板则负责操作输入解析、云台与发射机构控制、裁判系统交互以及机载计算平台通信。本节重点说明电控板侧的嵌入式软件设计；视觉定位和PX4飞行控制相关算法见前文视觉详细设计。

电控程序基于STM32CubeMX生成的HAL驱动和FreeRTOS实时操作系统开发，当前工程默认面向STM32F407，同时保留STM32F427的板级适配。工程采用C99和C++11混合编程，并按功能划分为四层：

\begin{itemize}
    \item \textbf{应用层（\texttt{app}）}：实现整机状态机、遥控/键鼠逻辑、云台、发射、自瞄接口、客户端UI和板间通信等功能；
    \item \textbf{基础服务层（\texttt{base}）}：封装电机、IMU、遥控器、裁判系统、视觉通信、CRC、PID、滤波器和连接状态监测等可复用模块；
    \item \textbf{接口层（\texttt{interface}）}：创建FreeRTOS任务，并将CAN、UART和USB中断回调分发到对应设备；
    \item \textbf{板级驱动层（\texttt{mcu}）}：包含STM32 HAL、FreeRTOS移植层以及CAN、UART、SPI、I\textsuperscript{2}C、定时器和独立看门狗等外设配置。
\end{itemize}

软件分层与主要数据流如\figref{fig:drone_embedded_software_architecture}所示。电控板不直接生成旋翼控制量，而是把飞手/云台手的方向指令或航点编号封装为\texttt{DRONE\_CONTROL}数据包发送给机载计算平台；机载计算平台完成安全约束和飞行控制接口转换后再交由PX4执行。该边界可避免云台与发射逻辑侵入飞行稳定回路，同时便于两部分独立调试。

\begin{figure}[hbtp!]
    \centering
    \begin{tikzpicture}[
        font=\small,
        layer/.style={draw, rounded corners=2pt, align=center, minimum width=14.0cm, minimum height=0.82cm, inner sep=4pt},
        ext/.style={draw, rounded corners=2pt, align=center, minimum width=3.3cm, minimum height=0.8cm, inner sep=3pt},
        flow/.style={-latex, thick},
        biflow/.style={latex-latex, thick}
    ]
        \node[ext] (operator) at (0,4.5) {遥控器/键鼠\\操作指令};
        \node[ext] (computer) at (5.2,4.5) {机载计算平台\\视觉、导航与飞行桥接};
        \node[ext] (referee) at (10.4,4.5) {裁判系统\\比赛状态与客户端UI};

        \node[layer] (app) at (5.2,3.0) {应用层：安全状态机、控制模式、云台、自瞄接口、发射机构、飞行指令桥接、UI};
        \node[layer] (base) at (5.2,1.6) {基础服务层：电机与IMU、遥控与裁判系统、视觉协议、PID/滤波、连接状态监测};
        \node[layer] (interface) at (5.2,0.2) {接口层：FreeRTOS任务调度、CAN/UART/USB中断回调与数据分发};
        \node[layer] (mcu) at (5.2,-1.2) {板级驱动层：STM32F407 HAL、CAN1/2、UART、SPI、I\textsuperscript{2}C、TIM、IWDG};

        \draw[flow] (operator) -- (app);
        \draw[biflow] (computer) -- (app);
        \draw[biflow] (referee) -- (app);
        \draw[biflow] (app) -- (base);
        \draw[biflow] (base) -- (interface);
        \draw[biflow] (interface) -- (mcu);
    \end{tikzpicture}
    \caption{空中机器人嵌入式软件分层与数据流}
    \label{fig:drone_embedded_software_architecture}
\end{figure}

\subsection{实时任务与中断设计}

程序采用“周期任务完成控制计算，中断回调完成数据接收”的组织方式。任务优先级按照姿态采样、电机闭环、整机控制、通信和调试依次降低，使\qty{1}{\milli\second}级控制链路优先于非关键数据显示。当前FreeRTOS任务配置如\tabref{tab:drone_freertos_tasks}所示。

\begin{table}[H]
    \centering
    \small
    \begin{tabular}{m{3.4cm} m{2.3cm} m{1.9cm} m{5.9cm}}
    \toprule
    任务 & CMSIS优先级 & 调度周期 & 主要职责 \\
    \midrule
    \texttt{imuTask}
    & Realtime & \qty{1}{\milli\second}
    & 读取板载IMU，更新姿态，并执行IMU加热闭环 \\
    \midrule
    \texttt{motorTask}
    & High & \qty{1}{\milli\second}
    & 更新电机反馈、位置/速度PID及最终电流控制量 \\
    \midrule
    \texttt{canTask}
    & High & 约\qty{1}{\milli\second}
    & 统一发送DJI电机与LK电机CAN控制帧，管理发送邮箱占用 \\
    \midrule
    \texttt{controlTask}
    & AboveNormal & \qty{1}{\milli\second}
    & 处理遥控输入、安全状态机、云台/发射目标及视觉工作模式 \\
    \midrule
    \texttt{minipcCommTask}
    & Normal & 分时调度
    & 接收机载计算平台数据，按\qty{5}{\milli\second}时隙发送通用、自瞄或飞行控制数据包 \\
    \midrule
    \texttt{refereeCommTask}
    & Normal & \qty{1}{\milli\second}
    & 解析裁判系统数据并刷新客户端UI \\
    \midrule
    \texttt{serialToolTask}
    & Low & \qty{20}{\milli\second}
    & 向Serial Studio或USB调试接口输出运行数据 \\
    \bottomrule
    \end{tabular}
    \caption{空中机器人FreeRTOS任务配置}
    \label{tab:drone_freertos_tasks}
\end{table}

CAN接收中断读取一帧标准帧后，依次交由电机驱动、超级电容通信和板间通信模块进行ID匹配；UART接收中断根据句柄将字节流分发给遥控器、机载计算平台、裁判系统或调试工具，遥控器同时使用UART空闲中断确定一帧结束。接收模块只在完整帧校验通过后更新数据并刷新连接时间戳，控制任务始终读取最近一帧有效数据，从而把耗时的控制计算留在任务上下文中。

\subsection{整机控制状态机}

电控软件使用\texttt{STOP}、\texttt{STARTUP}和\texttt{WORKING}三级状态机管理执行机构上电，状态转移如\figref{fig:drone_power_state_fsm}所示。遥控器失联、遥控器右拨杆置于下挡或任一云台电机离线均构成停止条件，系统无论处于启动还是工作状态都会立即返回\texttt{STOP}。

\begin{figure}[hbtp!]
    \centering
    \begin{tikzpicture}[
        font=\small,
        state/.style={draw, rounded corners=3pt, align=center, minimum width=3.0cm, minimum height=1.0cm, inner sep=4pt},
        flow/.style={-latex, thick}
    ]
        \node[state] (stop) at (0,0) {STOP\\停转并断电};
        \node[state] (startup) at (5.0,0) {STARTUP\\云台/拨盘回零};
        \node[state] (working) at (10.0,0) {WORKING\\正常控制};

        \draw[flow] (stop) -- node[midway, above=12pt, align=center] {遥控与云台链路正常\\且右拨杆非下挡} (startup);
        \draw[flow] (startup) -- node[midway, above=7pt] {初始化完成} (working);
        \draw[flow] (startup) to[bend left=28] node[below] {任一停止条件} (stop);
        \draw[flow] (working) to[bend left=42] node[below, align=center] {遥控失联/右拨杆下挡/云台离线} (stop);
    \end{tikzpicture}
    \caption{执行机构上电与安全状态机}
    \label{fig:drone_power_state_fsm}
\end{figure}

各状态的处理逻辑如下：

\begin{enumerate}
    \item \textbf{STOP}：关闭摩擦轮，并复位拨盘和控制模式。全部电机先进入\texttt{STOP}状态；当转速仍高于\qty{1200}{\degree\per\second}时执行闭环减速，低于阈值后再切换到\texttt{POWEROFF}并输出零电流，避免高速旋转时直接失去制动力矩。
    \item \textbf{STARTUP}：恢复电机输出并执行机械零点初始化。云台根据编码器零位慢速回正，回正阶段将位置环速度输出限制为\qty{180}{\degree\per\second}，Yaw和Pitch误差均小于\qty{5}{\degree}后解除限制；拨盘通过位置、电流、转速与超时联合判据完成回零。只有云台和拨盘均初始化完成后才进入工作状态。
    \item \textbf{WORKING}：根据遥控拨杆或键鼠输入选择手动瞄准、远距离自瞄、近距离自瞄、能量机关或基地模式，并周期更新云台、发射机构和机载计算平台数据包。
\end{enumerate}

独立看门狗与状态机形成双重保护。控制任务仅在遥控链路正常时刷新IWDG；一旦遥控数据超过\qty{100}{\milli\second}未更新，软件先切换至\texttt{STOP}并停止喂狗，若故障持续则由硬件看门狗复位控制器。

\subsection{云台与电机控制设计}

当前无人机执行机构的电机映射如\tabref{tab:drone_motor_mapping}所示。各电机对象只向上层暴露目标角度、目标速度或目标力矩接口，底层驱动统一完成编码器展开、连接检测、PID计算、输出限幅和CAN打包，降低应用逻辑与具体电机协议的耦合。

\begin{table}[hbtp!]
    \centering
    \begin{tabular}{m{3.0cm} m{3.1cm} m{3.2cm} m{4.3cm}}
    \toprule
    执行机构 & 电机型号 & 总线与ID & 控制方式 \\
    \midrule
    云台Yaw轴 & GM6020 & CAN1，ID 7 & 位置--速度串级闭环 \\
    \midrule
    云台Pitch轴 & MG4005（LK协议） & CAN2，ID 1 & 位置--速度串级闭环 \\
    \midrule
    左/右摩擦轮 & M2006各一台 & CAN1，ID 6/5 & 速度闭环，方向相反 \\
    \midrule
    拨盘 & M2006 & CAN1，ID 8 & 位置--速度串级闭环 \\
    \bottomrule
    \end{tabular}
    \caption{云台与发射机构电机映射}
    \label{tab:drone_motor_mapping}
\end{table}

云台位置--速度串级控制的计算形式为

\begin{equation}
    \omega_{\mathrm{ref}}
    =\operatorname{PID}_{p}(\theta_{\mathrm{ref}}-\theta_{\mathrm{fdb}})
    +\omega_{\mathrm{ff}},\qquad
    u=\operatorname{PID}_{s}(\omega_{\mathrm{ref}}-\omega_{\mathrm{fdb}})
    +u_{\mathrm{ff}},
    \label{equ:drone_gimbal_cascade_pid}
\end{equation}

其中$\theta$、$\omega$分别为角度和角速度，$u$为发送给电调的控制量，$\omega_{\mathrm{ff}}$来自视觉输出的目标角速度，$u_{\mathrm{ff}}$用于补偿机构重力与阻力。Pitch轴补偿量由Yaw、Pitch角度的多项式拟合得到，使不同姿态下的静态承载力矩直接进入速度环前馈，减小积分项负担。

云台支持三种反馈组合：正常工作时Yaw和Pitch均采用IMU姿态反馈；初始化时可切换为双编码器反馈；发射机构测试模式下Yaw采用编码器、Pitch采用IMU。目标角度在写入时即执行软件限位，当前Yaw范围为$[-80^\circ,80^\circ]$，Pitch范围为$[-24^\circ,40^\circ]$。手动、自瞄和能量机关模式分别使用独立PID参数组，模式切换时同步替换位置环与速度环参数，避免用单组参数兼顾不同动态需求。

板载IMU由独立的最高优先级任务以\qty{1}{\milli\second}周期更新。C型开发板配置使用BMI088进行角速度和加速度采样，软件完成安装方向变换、姿态解算和零偏补偿；加热电阻采用PWM和PID将传感器温度控制在约\qty{40}{\degreeCelsius}，减小环境温度变化对零偏的影响。

\subsection{自瞄接口与时间同步}

视觉算法向电控板返回Yaw/Pitch相对角度、目标角速度、距离、目标编号、发射标志和发射序号。电控板同时向视觉端反馈云台角度与角速度、补偿量、实测弹速、发射延迟、机器人ID和发射序号确认，从而形成“视觉解算目标---电控执行---结果回传”的闭环。

由于图像曝光、识别、串口传输和任务调度会引入延迟，直接把视觉相对角叠加到当前云台角度会造成动态误差。程序在\qty{1}{\milli\second}控制周期内保存最近64组云台角度，并按当前标定的\qty{22}{\milli\second}链路延迟回查历史姿态，目标角度按式\eqref{equ:drone_cv_time_alignment}计算：

\begin{equation}
    \boldsymbol{\theta}_{\mathrm{ref}}(t)
    =\boldsymbol{\theta}_{\mathrm{gimbal}}(t-\qty{22}{\milli\second})
    +\Delta\boldsymbol{\theta}_{\mathrm{cv}}(t)
    +\boldsymbol{\theta}_{\mathrm{offset}} .
    \label{equ:drone_cv_time_alignment}
\end{equation}

只有在自瞄使能、视觉目标包连接正常且视觉未返回空闲标志时，目标角度和角速度前馈才会写入云台控制器；自瞄数据超过\qty{25}{\milli\second}未刷新时，视觉控制和自动发射均被禁止。基地模式进入后的前\qty{500}{\milli\second}允许视觉快速拉回目标，随后清除视觉速度前馈并交还手动微调，兼顾快速对准和稳定操作。成功发射后，电控板将收到的\texttt{shoot\_id}原样回传，供视觉端确认该次发射指令已执行。

\subsection{发射机构软件设计}

发射机构由双摩擦轮和拨盘构成。摩擦轮打开后，程序等待两侧电机连接正常且转速接近目标值，再置位摩擦轮就绪标志；拨盘每次发射将目标角度推进\qty{45}{\degree}。当前默认最小发射间隔为\qty{40}{\milli\second}，单次发射许可由下式统一门控：

\begin{equation}
    S_{\mathrm{fire}}=S_{\mathrm{fric\_ready}}
    \land \neg S_{\mathrm{blocked}}
    \land S_{\mathrm{heat\_ok}}
    \land (t-t_{\mathrm{last}}>t_{\mathrm{CD}}).
    \label{equ:drone_shoot_interlock}
\end{equation}

裁判系统在线时，软件实时读取弹速上限、枪口热量上限、冷却速率和实测初速度。为弥补裁判系统热量数据在连发时的更新延迟，程序在本地按每发10点热量累加、按冷却速率每毫秒递减，并保留大于2.5发弹丸的热量余量后才允许下一次发射。实测弹速经过三点均值滤波后与\qty{24}{\meter\per\second}目标值比较，按误差小步调整摩擦轮目标转速，以降低个体弹丸和电池电压变化带来的初速偏差。

拨盘防卡弹逻辑同时监测位置误差和电机电流。当大位置误差与反向大电流持续出现时，程序暂停发射、记录当前拨盘位置并执行短时反转卸载，随后回到上一弹位继续工作。摩擦轮转速突变还用于估计从下发拨盘指令到弹丸实际离膛的机械延迟，该结果通过视觉数据包反馈，可用于后续目标预测补偿。

\subsection{飞行指令桥接与通信协议}

电控板从DR16遥控器获得拨杆、摇杆、鼠标和键盘状态。云台手可在手动控制、自瞄、能量机关和基地等模式之间切换；在允许控制飞机时，键盘输入被转换为两类互斥的飞行指令：

\begin{itemize}
    \item \textbf{方向控制（\texttt{DIRECTION}）}：以位标志同时表示前、后、左、右、左偏航和右偏航，供机载计算平台生成连续速度指令；
    \item \textbf{航点控制（\texttt{SETPOINT}）}：发送0至3号预设航点索引，供导航节点生成位置目标；
    \item \textbf{无效状态（\texttt{INVALID}）}：每个控制周期开始时默认写入无效指令；只有当前操作模式明确授权后才改写为方向或航点控制，防止旧指令残留。
\end{itemize}

电控板与机载计算平台通过UART通信，串口配置为115200 bit/s、8位数据位、1位停止位、无校验和无流控。协议帧结构为

\begin{equation}
    \underbrace{\texttt{0x23}}_{\mathrm{SOF}}
    +\underbrace{\texttt{pack\_id}}_{\text{方向与类型}}
    +\underbrace{\texttt{data\_len}}_{\text{载荷长度}}
    +\underbrace{\texttt{payload}}_{\text{业务数据}}
    +\underbrace{\texttt{CRC16}}_{\text{整帧校验}} .
    \label{equ:drone_cv_uart_frame}
\end{equation}

其中\texttt{pack\_id}高四位表示传输方向，低四位表示通用、自瞄、导航、比赛信息、无人机控制或无人机状态数据。无人机控制包固定为3字节载荷，依次包含控制模式、方向位标志和航点索引；机载计算平台反向发送飞控连接、位姿可用性、导航可用性及当前控制源，电控板将这些状态显示在客户端UI上。主要链路设计见\tabref{tab:drone_communication_interfaces}。

\begin{table}[hbtp!]
    \centering
    \begin{tabular}{m{2.4cm} m{2.8cm} m{4.0cm} m{4.6cm}}
    \toprule
    链路 & 物理接口 & 主要数据 & 异常处理 \\
    \midrule
    遥控器 & UART3，DR16 & 拨杆、摇杆、键盘、鼠标 & \qty{100}{\milli\second}未刷新即进入\texttt{STOP} \\
    \midrule
    机载计算平台 & UART1，115200 bit/s & 自瞄、方向/航点、飞行状态 & CRC16与长度校验；按数据类型设置\qty{25}{\milli\second}至\qty{300}{\milli\second}超时 \\
    \midrule
    裁判系统 & UART6，115200 bit/s & 比赛状态、热量、弹速、客户端UI & \qty{1000}{\milli\second}超时；离线时使用预设发射参数 \\
    \midrule
    电机与板间通信 & CAN1/CAN2 & 电机反馈/控制、板间状态 & 电机\qty{50}{\milli\second}离线后控制量清零；按标准帧ID分发 \\
    \midrule
    调试接口 & UART或USB CDC & 关键状态与调参数据 & 低优先级、可编译关闭，不占用控制链路 \\
    \bottomrule
    \end{tabular}
    \caption{空中机器人主要通信接口}
    \label{tab:drone_communication_interfaces}
\end{table}

裁判系统客户端UI用于显示当前自瞄模式、飞行控制模式、六个方向位、0至3号航点、飞行状态以及摩擦轮是否就绪。该界面把电控板内部状态直接反馈给操作手，可在比赛中及时发现飞控未连接、导航不可用或控制源异常等问题。

\subsection{故障检测与降级策略}

所有关键通信对象统一使用时间戳式连接监测：接收中断在有效帧到达时刷新时间戳，周期任务比较当前时刻与最后更新时间，超过各自阈值后将链路标记为离线。针对不同故障，软件采用分级降级策略：

\begin{enumerate}
    \item 遥控器或云台关键电机离线时，整机立即进入\texttt{STOP}，关闭发射并停止所有电机输出；
    \item 单个普通电机超过\qty{50}{\milli\second}无反馈时，该电机控制量独立清零，避免失联后保持旧电流；
    \item 自瞄数据超时、CRC错误、长度错误或视觉返回空闲状态时，不更新视觉目标并禁止自动发射，但保留人工接管能力；
    \item 飞行控制命令默认置为\texttt{INVALID}，方向控制和航点控制不能同时生效，降低粘连按键或历史数据造成持续飞行的风险；
    \item 裁判系统离线时不再信任其热量与弹速反馈，发射模块退回预设参数；调试串口与UI故障不影响实时控制任务。
\end{enumerate}

软件联调应覆盖正常启动、遥控失联、视觉数据中断、CAN电机掉线、裁判系统掉线、拨盘卡弹和看门狗复位等故障注入场景，并确认所有情况下均能进入预期安全状态。飞行指令联调时，应同时记录电控板发送的\texttt{DRONE\_CONTROL}、机载计算平台返回的\texttt{DRONE\_STATUS}及PX4实际控制源，验证方向位、航点编号和控制权切换的一致性。

\subsection{构建、烧录与调试}

工程使用CMake组织\texttt{base}、\texttt{app}和\texttt{interface}静态库，采用\texttt{arm-none-eabi-gcc/g++}交叉编译并生成ELF、HEX和BIN三种产物。调试构建使用\texttt{-Og -g}保留符号，发布构建使用\texttt{-Ofast}；通过根目录的MCU选项可在STM32F407和STM32F427板级工程之间切换。硬件外设变更由STM32CubeMX维护，在线调试可使用Keil或OpenOCD，运行数据通过Serial Studio、USB CDC和裁判系统客户端UI观察。

为降低参数修改对安全逻辑的影响，调试过程按“单模块闭环---整机状态机---视觉通信---飞行联调”的顺序进行。云台角度限位、PID输出上限、摩擦轮转速、发射热量余量以及各链路超时阈值均应在版本管理中留存，并在正式比赛固件发布前执行完整编译、烧录校验和故障注入回归测试。

\clearpage
